\begin{proof}[Proof of \cref{obj:lem:measurable-component-contraction}]
\begin{enumerate}
\item \textit{Conditional factorization.}
Fix a covariate vector \(x\in[0,1]^n\). To identify the signs used by each record, write the envelope and micro frames in \cref{obj:def:cosine-family} as
\[
g_{\mathrm L}(v)=
\begin{cases}
\cos^2\!\left(\dfrac{\pi(v-x_0)}{2h_{\mathrm L}}\right),
& |v-x_0|\le h_{\mathrm L},\\
0,& |v-x_0|>h_{\mathrm L},
\end{cases}
\qquad
\phi_{\mathrm L,\eta}(v)=
\begin{cases}
\cos\!\left(\dfrac{\pi(v-\eta\delta_{\mathrm L})}{2\delta_{\mathrm L}}\right),
& |v-\eta\delta_{\mathrm L}|\le\delta_{\mathrm L},\\
0,& |v-\eta\delta_{\mathrm L}|>\delta_{\mathrm L},
\end{cases}
\]
for \(v\in[0,1]\) and \(\eta\in\mathbb Z\), where \(x_0=1/2\). Thus
\[
q(v)=g_{\mathrm L}(v)^2,
\qquad
S_\lambda(v)=g_{\mathrm L}(v)
 \sum_{\eta\in\mathcal J_{\mathrm L}}(2\lambda_\eta-1)\phi_{\mathrm L,\eta}(v),
\]
with the finite active index set
\[
\mathcal J_{\mathrm L}
=
\left\{
\eta\in\mathbb Z:
\begin{array}{l}
\displaystyle
\left\lfloor\frac{x_0-h_{\mathrm L}}{\delta_{\mathrm L}}\right\rfloor-1
\le \eta\le
\left\lceil\frac{x_0+h_{\mathrm L}}{\delta_{\mathrm L}}\right\rceil+1,\\
\eta\delta_{\mathrm L}-\delta_{\mathrm L}<x_0+h_{\mathrm L},
\quad
x_0-h_{\mathrm L}<\eta\delta_{\mathrm L}+\delta_{\mathrm L}
\end{array}
\right\}.
\]
Define the record and component sign sets by
\[
J_i(x)=
 \{\eta\in\mathcal J_{\mathrm L}:
   g_{\mathrm L}(x_i)\phi_{\mathrm L,\eta}(x_i)\ne0\},
\qquad
J_C(x)=\bigcup_{i\in C}J_i(x).
\]
The graph and its ordered component list are therefore specified by
\[
i\sim_{G_x}i'
\quad\Longleftrightarrow\quad
i\ne i'\ \text{ and }\ J_i(x)\cap J_{i'}(x)\ne\varnothing,
\]
\[
\mathscr C(x)
=
\bigl(C:\ C\text{ is a connected component of }G_x\bigr),
\qquad
\text{listed by increasing }\min C.
\]
For \(n=0\), this list is empty.

We use binary mark configurations and their signed likelihood coordinates through the encoding
\[
z_i=(A_i,Y_i)\in\{0,1\}^2,
\qquad
U_i(z)=2A_i-1,
\qquad
V_{Y,i}(z)=2Y_i-1.
\]
This identifies binary configurations with the signed configurations in the statement. The likelihood expansion supplied by
\cref{obj:lem:positive-family-certificate} shows that the conditional likelihood at record \(i\) depends on \(\lambda\) only through \(J_i(x)\). Its quadratic term uses the same sign set.

Distinct components have disjoint sign sets:
\[
C,C'\in\mathscr C(x),\ C\ne C'
\quad\Longrightarrow\quad
J_C(x)\cap J_{C'}(x)=\varnothing.
\]
Indeed, a common sign would be used by a vertex in each component and would join them by an edge. Since the signs are independent and fair, summing separately over these disjoint sign blocks gives
\[
E_\lambda\prod_{i=1}^n
 \frac{L_\lambda(x_i,U_i(z),V_{Y,i}(z))}{4}
=
\prod_{C\in\mathscr C(x)}
 E_\lambda\prod_{i\in C}
 \frac{L_\lambda(x_i,U_i(z),V_{Y,i}(z))}{4}
=
\prod_{C\in\mathscr C(x)}Q_C(\{z|_C\}).
\]
Unused signs average to one. The component vertex sets partition the record indices, so
\[
4^{-n}
=
\prod_{C\in\mathscr C(x)}4^{-|C|}
=
\prod_{C\in\mathscr C(x)}P_{0,C}(\{z|_C\}).
\]
Both identities also hold for \(n=0\), with empty products equal to one.
% lean: CausalSmith.Stat.PrivateCateRoughdesign.fullAlternativeMass_component_factorization
% lean: CausalSmith.Stat.PrivateCateRoughdesign.nullMass_component_factorization

\item \textit{Singleton cancellation.}
We first establish the pointwise sign identities needed below. Define, for \(v\in[0,1]\),
\[
\eta_v=\left\lfloor\frac{v}{\delta_{\mathrm L}}\right\rfloor,
\qquad
\omega_{\mathrm L}(v)
=\frac{\pi(v-\eta_v\delta_{\mathrm L})}{2\delta_{\mathrm L}}.
\]
Then \(0\le\omega_{\mathrm L}(v)<\pi/2\), and the only possibly nonzero frames are the two neighboring ones:
\[
\phi_{\mathrm L,\eta_v}(v)=\cos\omega_{\mathrm L}(v),
\qquad
\phi_{\mathrm L,\eta_v+1}(v)=\sin\omega_{\mathrm L}(v),
\qquad
\phi_{\mathrm L,\eta}(v)=0
\quad(\eta\notin\{\eta_v,\eta_v+1\}).
\]
Consequently their squared values sum to one. When \(q(v)>0\), every nonzero frame at \(v\) belongs to the active set; when \(q(v)=0\), multiplying by \(q(v)\) annihilates every term. Hence, including support boundaries,
\[
q(v)\sum_{\eta\in\mathcal J_{\mathrm L}}
 \phi_{\mathrm L,\eta}(v)^2=q(v).
\]
Fair independent signs satisfy
\[
E_\lambda(2\lambda_\eta-1)=0,
\qquad
E_\lambda\bigl[(2\lambda_\eta-1)(2\lambda_{\eta'}-1)\bigr]
=\mathbf 1\{\eta=\eta'\}.
\]
Expanding the perturbation and its square therefore yields
\[
E_\lambda S_\lambda(v)=0,
\qquad
E_\lambda S_\lambda(v)^2
=
q(v)\sum_{\eta\in\mathcal J_{\mathrm L}}\phi_{\mathrm L,\eta}(v)^2
=q(v).
\]
Substitution into the likelihood expansion in
\cref{obj:lem:positive-family-certificate} gives the pointwise identity
\[
E_\lambda L_\lambda(v,U,V_Y)=1
\qquad
(v\in[0,1],\ U,V_Y\in\{-1,1\}).
\]
Thus, if \(C=\{i\}\), every binary mark pair has mass
\[
Q_C(\{z\})
=\frac14E_\lambda L_\lambda(x_i,U_i(z),V_{Y,i}(z))
=\frac14
=P_{0,C}(\{z\}),
\]
and \(Q_C=P_{0,C}\).
% lean: CausalSmith.Stat.PrivateCateRoughdesign.conditionalLikelihood_sign_sum_exact
% lean: CausalSmith.Stat.PrivateCateRoughdesign.alternativeComponentLaw_eq_null_of_card_one

\item \textit{Component curvature.}
Fix a component \(C\), and put \(m=|C|\). The two-frame property and Cauchy--Schwarz imply, for every \(\lambda\) and \(v\),
\[
S_\lambda(v)^2
\le
2q(v)\sum_{j\in\mathcal J_{\mathrm L}}\phi_{\mathrm L,j}(v)^2
=2q(v).
\]
Since \(0\le q(v)\le1\), this gives the coefficient bounds
\[
|S_\lambda(v)|\le\frac32,
\qquad
|q(v)-S_\lambda(v)^2|\le q(v)\le1.
\]

For \(i\in C\), \(\lambda\), \(z_i\in\{0,1\}^2\), and every real amplitude \(u\), define the polynomial mark mass
\[
\begin{aligned}
p_{i,\lambda}(z_i;u)
=\frac14\bigl[
1
&+U_i(z)uS_\lambda(x_i)\\
&+V_{Y,i}(z)u(u^2q(x_i)-1)S_\lambda(x_i)\\
&+U_i(z)V_{Y,i}(z)u^2
   \bigl(q(x_i)-S_\lambda(x_i)^2\bigr)
\bigr].
\end{aligned}
\]
Its factorization is
\[
p_{i,\lambda}(z_i;u)
=
\frac{1+U_i(z)uS_\lambda(x_i)}2
\frac{1+V_{Y,i}(z)
 \bigl(-uS_\lambda(x_i)+U_i(z)u^2q(x_i)\bigr)}2.
\]
For \(0\le u\le1/10\), the two perturbations satisfy
\[
|uS_\lambda(x_i)|\le\frac3{20},
\qquad
\bigl|-uS_\lambda(x_i)+U_i(z)u^2q(x_i)\bigr|
\le\frac3{20}+\frac1{100}
=\frac4{25}.
\]
The masses are therefore positive. Summing first over the outcome mark and then over the treatment mark gives
\[
\sum_{z_i\in\{0,1\}^2}p_{i,\lambda}(z_i;u)=1,
\qquad
\sum_{z_i\in\{0,1\}^2}|p_{i,\lambda}(z_i;u)|=1.
\]

Differentiation gives
\[
\begin{aligned}
\frac{\partial p_{i,\lambda}}{\partial u}(z_i;u)
=\frac14\bigl[
&U_i(z)S_\lambda(x_i)
+V_{Y,i}(z)(3u^2q(x_i)-1)S_\lambda(x_i)\\
&+2U_i(z)V_{Y,i}(z)u
 \bigl(q(x_i)-S_\lambda(x_i)^2\bigr)
\bigr],
\end{aligned}
\]
\[
\frac{\partial^2p_{i,\lambda}}{\partial u^2}(z_i;u)
=\frac14\bigl[
6V_{Y,i}(z)u q(x_i)S_\lambda(x_i)
+2U_i(z)V_{Y,i}(z)
 \bigl(q(x_i)-S_\lambda(x_i)^2\bigr)
\bigr].
\]
On \(0\le u\le1/10\), we have
\[
|3u^2q(x_i)-1|\le1.
\]
Taking absolute values of the displayed Fourier coefficients and summing over the four mark pairs consequently gives
\[
\begin{aligned}
\sum_{z_i}
 \left|\frac{\partial p_{i,\lambda}}{\partial u}(z_i;u)\right|
&\le
|S_\lambda(x_i)|
+|3u^2q(x_i)-1|\,|S_\lambda(x_i)|\\
&\qquad
+2u|q(x_i)-S_\lambda(x_i)^2|\\
&\le \frac32+\frac32+\frac15
\le4,
\end{aligned}
\]
and
\[
\begin{aligned}
\sum_{z_i}
 \left|\frac{\partial^2p_{i,\lambda}}{\partial u^2}(z_i;u)\right|
&\le
6u q(x_i)|S_\lambda(x_i)|
+2|q(x_i)-S_\lambda(x_i)^2|\\
&\le\frac9{10}+2
\le4.
\end{aligned}
\]

Define the averaged component array, for every \(u\in\mathbb R\), by
\[
F_{C,z}(u)=E_\lambda\prod_{i\in C}p_{i,\lambda}(z_i;u),
\qquad
z\in(\{0,1\}^2)^C.
\]
The second product derivative consists of one term for each twice-differentiated factor and one term for each ordered pair of distinct once-differentiated factors:
\[
\begin{aligned}
\frac{d^2}{du^2}\prod_{i\in C}p_{i,\lambda}(z_i;u)
={}&
\sum_{i\in C}
 \frac{\partial^2p_{i,\lambda}}{\partial u^2}(z_i;u)
 \prod_{i'\in C\setminus\{i\}}p_{i',\lambda}(z_{i'};u)\\
&+
\sum_{\substack{i,i'\in C\\i\ne i'}}
 \frac{\partial p_{i,\lambda}}{\partial u}(z_i;u)
 \frac{\partial p_{i',\lambda}}{\partial u}(z_{i'};u)
 \prod_{i''\in C\setminus\{i,i'\}}p_{i'',\lambda}(z_{i''};u).
\end{aligned}
\]
For finite arrays, summing the absolute value of a tensor product factors into the product of the absolute sums. Applying the preceding bounds and then averaging over signs yields
\[
\sum_z|F_{C,z}''(u)|
\le4m+16m(m-1)
\le16m^2
\qquad(0\le u\le1/10).
\]
At zero amplitude all factors equal \(1/4\), while
\[
\frac{\partial p_{i,\lambda}}{\partial u}(z_i;0)
=
\frac{U_i(z)-V_{Y,i}(z)}4S_\lambda(x_i).
\]
The pointwise sign cancellation from the preceding step thus gives
\[
F_{C,z}(0)=4^{-m},
\qquad
F_{C,z}'(0)=0.
\]
For an empty vertex set, the product is identically one and both derivatives vanish, so the same conclusions and curvature bound apply.
% lean: CausalSmith.Stat.PrivateCateRoughdesign.amplitudeMark_coefficients_valid
% lean: CausalSmith.Stat.PrivateCateRoughdesign.amplitudeMarkSlope_abs_sum_le_four
% lean: CausalSmith.Stat.PrivateCateRoughdesign.amplitudeMarkCurvature_abs_sum_le_four
% lean: CausalSmith.Stat.PrivateCateRoughdesign.componentAmplitudeCurvature_abs_sum_le
% lean: CausalSmith.Stat.PrivateCateRoughdesign.componentAmplitudeSlope_zero

\item \textit{Quadratic remainder and common overlap.}
The calibration implies
\[
0<t\le\frac1{16384},
\qquad
0\le\sqrt t\le\frac1{10}.
\]
At this amplitude, the likelihood expansion in
\cref{obj:lem:positive-family-certificate} identifies the component arrays as
\[
a_z:=Q_C(\{z\})=F_{C,z}(\sqrt t),
\qquad
b_z:=P_{0,C}(\{z\})=F_{C,z}(0)=4^{-m}.
\]
To control their entire absolute difference with one scalar remainder, define
\[
\sigma_C(z)=
\begin{cases}
1,&a_z-b_z\ge0,\\
-1,&a_z-b_z<0,
\end{cases}
\qquad
\mathcal F_C(u)=\sum_z\sigma_C(z)F_{C,z}(u)
\quad(u\in\mathbb R).
\]
Then
\[
\mathcal F_C'(0)=0,
\qquad
|\mathcal F_C''(u)|
\le\sum_z|F_{C,z}''(u)|
\le16m^2
\quad(0\le u\le\sqrt t).
\]
The mean-value inequality gives
\[
|\mathcal F_C'(u)|\le16m^2u
\qquad(0\le u\le\sqrt t).
\]
Integrating the first derivative, and using the defining signs, proves
\[
\begin{aligned}
\sum_z|a_z-b_z|
&=\mathcal F_C(\sqrt t)-\mathcal F_C(0)\\
&\le\int_0^{\sqrt t}|\mathcal F_C'(u)|\,du\\
&\le\int_0^{\sqrt t}16m^2u\,du
=8tm^2.
\end{aligned}
\]

Define the common entries and total overlap by
\[
c_z:=\min(a_z,b_z),
\qquad
\alpha:=\sum_zc_z.
\]
Both arrays are nonnegative and sum to one, by the mark-mass normalization above and the \(4^m\) null configurations. Therefore
\[
0\le\alpha\le1,
\qquad
\sum_z(a_z-c_z)=\sum_z(b_z-c_z)=1-\alpha,
\]
and
\[
\sum_z|a_z-b_z|
=\sum_z(a_z+b_z-2c_z)
=2(1-\alpha).
\]
For every mark event
\(\mathcal E_C\subseteq(\{0,1\}^2)^C\), each residual sum lies in \([0,1-\alpha]\), so
\[
\left|
 \sum_{z\in\mathcal E_C}a_z-\sum_{z\in\mathcal E_C}b_z
\right|
=
\left|
 \sum_{z\in\mathcal E_C}(a_z-c_z)
 -
 \sum_{z\in\mathcal E_C}(b_z-c_z)
\right|
\le1-\alpha.
\]
Taking the supremum over events gives
\[
d_{\mathrm{TV}}(Q_C,P_{0,C})
\le1-\alpha
=\frac12\sum_z|a_z-b_z|
\le4tm^2.
\]
In fact, the event
\[
\mathcal E_C^+:=\{z:b_z\le a_z\}
\]
has probability difference
\[
Q_C(\mathcal E_C^+)-P_{0,C}(\mathcal E_C^+)
=\sum_z(a_z-c_z)=1-\alpha.
\]
Thus the total variation distance equals \(1-\alpha\), as will also describe the coupling's failure mass.
% lean: CausalSmith.Stat.PrivateCateRoughdesign.finite_array_second_derivative_remainder_le
% lean: CausalSmith.Stat.PrivateCateRoughdesign.alternativeMass_abs_sub_nullMass_sum_le
% lean: CausalSmith.Stat.PrivateCateRoughdesign.overlapMass_deficit_le
% lean: CausalSmith.Stat.PrivateCateRoughdesign.component_TV_le_quadratic
% lean: CausalSmith.Stat.PrivateCateRoughdesign.component_TV_eq_overlap_deficit

\item \textit{Component coupling.}
Use the masses \(\Gamma_C(z,w)\) defined in
\cref{obj:def:common-mass-coupling}, with the arrays just specified. When \(\alpha<1\), summing its formula over the second coordinate gives
\[
\sum_w\Gamma_C(z,w)
=
c_z+\frac{(a_z-c_z)\sum_w(b_w-c_w)}{1-\alpha}
=a_z.
\]
Similarly,
\[
\sum_z\Gamma_C(z,w)
=
c_w+\frac{\sum_z(a_z-c_z)(b_w-c_w)}{1-\alpha}
=b_w.
\]
When \(\alpha=1\), all nonnegative residual entries have zero sum and hence vanish individually:
\[
a_z=c_z=b_z
\qquad\text{for every }z.
\]
The diagonal branch then has the same row and column sums. In both cases the masses are nonnegative and sum to one, establishing that \(\Gamma_C\) couples \(Q_C\) and \(P_{0,C}\).

At every configuration, at least one residual entry vanishes:
\[
(a_z-c_z)(b_z-c_z)=0.
\]
Consequently the diagonal and off-diagonal masses are
\[
\Gamma_C(z,z)=c_z,
\qquad
\sum_{z\ne w}\Gamma_C(z,w)=1-\alpha.
\]
For a singleton component, the cancellation already proved gives \(\alpha=1\).
% lean: CausalSmith.Stat.PrivateCateRoughdesign.componentCoupling_isCoupling
% lean: CausalSmith.Stat.PrivateCateRoughdesign.componentCouplingMass_diagonal
% lean: CausalSmith.Stat.PrivateCateRoughdesign.overlapMass_eq_one_of_card_one

\item \textit{Measurable dataset coupling and its marginals.}
For full binary mark vectors \(z,w\in(\{0,1\}^2)^n\), define
\[
p_{\mathrm{joint}}(x,z,w)
:=
\prod_{C\in\mathscr C(x)}
 \Gamma_C(z|_C,w|_C),
\qquad
D_{\mathrm{attach}}(x,z)
:=
\bigl((x_i,A_i,Y_i)\bigr)_{i=1}^n,
\]
where \(z_i=(A_i,Y_i)\). The conditional dataset measure is
\[
\gamma_x
:=
\sum_{z,w}p_{\mathrm{joint}}(x,z,w)\,
\delta_{\left(D_{\mathrm{attach}}(x,z),
               D_{\mathrm{attach}}(x,w)\right)}.
\]

Each envelope and frame is continuous, including at its support endpoints. The edge predicates defining \(G_x\) are finite unions of Borel nonzero predicates. For every fixed labeled graph \(\mathsf G\) on the record indices, the set
\[
\mathcal X_{\mathsf G}:=\{x\in[0,1]^n:G_x=\mathsf G\}
\]
is therefore Borel. There are finitely many such graphs. On each \(\mathcal X_{\mathsf G}\), the ordered component list is fixed, and the component masses are finite sums of products of Borel functions. Taking minima, finite sums, and the two overlap branches preserves measurability; the denominator in the branch \(\alpha<1\) is positive. It follows that \(x\mapsto p_{\mathrm{joint}}(x,z,w)\) is Borel. The attachment maps are Borel as well, so the finite atomic formula makes \(x\mapsto\gamma_x\) measurable.

The component marginal identities and factorization give, for every \(x,z,w\),
\[
\begin{aligned}
\sum_w p_{\mathrm{joint}}(x,z,w)
&=\prod_{C\in\mathscr C(x)}Q_C(\{z|_C\})\\
&=E_\lambda\prod_{i=1}^n
 \frac{L_\lambda(x_i,U_i(z),V_{Y,i}(z))}{4},
\end{aligned}
\qquad
\sum_zp_{\mathrm{joint}}(x,z,w)=4^{-n}.
\]
In particular, \(\gamma_x\) is a probability measure.

Define the shared covariate law and integrated coupling by
\[
\nu_X:=\bigotimes_{i=1}^n\operatorname{Unif}[0,1],
\qquad
\gamma(\mathcal E)
:=
\int_{[0,1]^n}\gamma_x(\mathcal E)\,\nu_X(dx)
\]
for every measurable dataset-pair event \(\mathcal E\).
The alternatives and null have this common covariate law. Integrating the first conditional marginal and interchanging the finite sign average with the integral yields the uniform mixture of alternative product laws; integrating the second yields the fair-null product law. Explicitly, with
\[
\operatorname{pr}_1(D,D')=D,
\qquad
\operatorname{pr}_2(D,D')=D',
\]
we obtain
\[
(\operatorname{pr}_1)_\#\gamma=Q_n,
\qquad
(\operatorname{pr}_2)_\#\gamma=(P_0^{\mathrm{obs}})^{\otimes n},
\qquad
\gamma(\mathcal O^n\times\mathcal O^n)=1.
\]
For \(n=0\), the mark space and dataset space each have one element; the empty component product gives the unit mass on the pair of empty datasets.
% lean: CausalSmith.Stat.PrivateCateRoughdesign.measurable_conditionalDatasetCoupling
% lean: CausalSmith.Stat.PrivateCateRoughdesign.commonMassCoupling_isCoupling

\item \textit{Component cost and tree probabilities.}
Assume now that \(n\delta_{\mathrm L}\le1/128\), and first consider \(n\ge1\). For component configurations define
\[
d_{\mathrm H,C}(z,w)
:=\sum_{i\in C}\mathbf 1\{z_i\ne w_i\},
\qquad
c_{\mathrm H,C}(x)
:=\sum_{z,w}d_{\mathrm H,C}(z,w)\Gamma_C(z,w).
\]
Because \(d_{\mathrm H,C}\) is zero on the diagonal and at most \(m\) elsewhere,
\[
c_{\mathrm H,C}(x)
\le m(1-\alpha)
\le4tm^3.
\]
For \(m=1\), the coupling is diagonal and its cost is zero. The shared covariates ensure that dataset Hamming distance is exactly the sum of component mark distances. Thus
\[
\begin{aligned}
\int d_{\mathrm H}(D,D')\,\gamma(dD,dD')
&=\int\sum_{C\in\mathscr C(x)}c_{\mathrm H,C}(x)\,\nu_X(dx)\\
&\le4t\sum_{m=2}^n m^3\int N_m(x)\,\nu_X(dx),
\end{aligned}
\]
where, for every integer \(m\ge0\), the component count is
\[
N_m(x)
:=\#\{C\in\mathscr C(x):|C|=m\}.
\]
In particular \(N_0(x)=0\), and \(N_m(x)=0\) for \(m>n\).

To bound these counts, a shared-sign edge has the geometric restrictions
\[
i\sim_{G_x}i'
\quad\Longrightarrow\quad
|x_i-x_0|<h_{\mathrm L},
\quad
|x_{i'}-x_0|<h_{\mathrm L},
\quad
|x_i-x_{i'}|<2\delta_{\mathrm L}.
\]
Indeed, the common nonzero weighted frame places both covariates in the open macro support and within \(\delta_{\mathrm L}\) of the same frame center. The last inequality follows by the triangle inequality.

Fix \(2\le m\le n\), an \(m\)-element record set \(C\), and a bijection
\[
e_C:\{1,\ldots,m\}\longrightarrow C.
\]
Let
\[
\mathscr T_m
:=\{\text{undirected trees on the labeled vertex set }\{1,\ldots,m\}\}.
\]
For \(\mathsf T\in\mathscr T_m\), define its shared-edge event by
\[
\mathcal E_{\mathsf T,C}
:=
\left\{
x:
e_C(i)\sim_{G_x}e_C(i')
\ \text{for every edge }\{i,i'\}\text{ of }\mathsf T
\right\}.
\]
A connected tree admits a parent-before-child order: start at a root and repeatedly adjoin a vertex adjacent to the vertices already selected. Connectivity permits this until all vertices are selected. Write this order and its parent positions as
\[
\rho_{\mathsf T}:\{1,\ldots,m\}\longrightarrow\{1,\ldots,m\}
\quad\text{a permutation},
\qquad
\pi_{\mathsf T}(i)\in\{1,\ldots,i-1\}
\quad(2\le i\le m),
\]
with
\[
\{\rho_{\mathsf T}(i),\rho_{\mathsf T}(\pi_{\mathsf T}(i))\}
\text{ an edge of }\mathsf T.
\]
Since \(m\ge2\), the first edge localizes the root. The shared-edge event is contained in the region
\[
\begin{split}
\mathcal R_{\mathsf T,C}:=\bigl\{x:\ &
|x_{e_C(\rho_{\mathsf T}(1))}-x_0|<h_{\mathrm L},\\
&
|x_{e_C(\rho_{\mathsf T}(i))}
 -x_{e_C(\rho_{\mathsf T}(\pi_{\mathsf T}(i)))}
 |<2\delta_{\mathrm L}
\quad(2\le i\le m)\bigr\}.
\end{split}
\]
The root interval has length at most \(2h_{\mathrm L}\), and, conditional on the preceding coordinates, each child interval has length at most \(4\delta_{\mathrm L}\). Integrating the last child first and repeating gives
\[
\nu_X(\mathcal E_{\mathsf T,C})
\le\nu_X(\mathcal R_{\mathsf T,C})
\le2h_{\mathrm L}(4\delta_{\mathrm L})^{m-1}.
\]
Coordinates outside \(C\) integrate to one.

If \(C\) is a graph component, its connected induced graph contains a spanning tree. Hence
\[
\{x:C\in\mathscr C(x)\}
\subseteq
\bigcup_{\mathsf T\in\mathscr T_m}\mathcal E_{\mathsf T,C}.
\]
The labeled-tree count is
\[
|\mathscr T_m|=m^{m-2},
\]
by \citep[Section 5.6, Theorem 5.40]{KellerTrotterAppliedCombinatorics}. A union bound, followed by summation over the \(\binom nm\) record subsets, proves
\[
\int N_m(x)\,\nu_X(dx)
\le
\binom nm\,2h_{\mathrm L}m^{m-2}
 (4\delta_{\mathrm L})^{m-1}.
\]
All geometric restrictions used here are strict nonzero-support restrictions, so the estimate includes every support-boundary configuration.
% lean: CausalSmith.Stat.PrivateCateRoughdesign.commonMassCoupling_hamming_cost_le_cubic_sum
% lean: CausalSmith.Stat.PrivateCateRoughdesign.uniform_labeled_tree_probability_bound
% lean: CausalSmith.Stat.PrivateCateRoughdesign.lintegral_sizeComponentCount_le_of_uniform_labeled_tree_bound

\item \textit{Summation of the transport envelope.}
Using \(\binom nm\le n^m/m!\) in the preceding cost bound gives
\[
\begin{aligned}
\int d_{\mathrm H}(D,D')\,\gamma(dD,dD')
&\le
\sum_{m=2}^n
4tm^3\binom nm\,2h_{\mathrm L}m^{m-2}
 (4\delta_{\mathrm L})^{m-1}\\
&\le
32tn^2h_{\mathrm L}\delta_{\mathrm L}
\sum_{m=2}^n
\frac{m^{m+1}}{m!}(4n\delta_{\mathrm L})^{m-2}.
\end{aligned}
\]
For \(n=1\), both sums are empty.

The nonnegative exponential series gives
\[
\frac{m^m}{m!}
\le
\sum_{j_{\mathrm{exp}}=0}^{\infty}
 \frac{m^{j_{\mathrm{exp}}}}{j_{\mathrm{exp}}!}
=\exp(m)
=\exp(1)^m.
\]
Define the series argument
\[
\zeta_{\mathrm{tree}}
:=4\exp(1)n\delta_{\mathrm L}.
\]
The sparsity cap and \(\exp(1)<3\) imply
\[
0\le\zeta_{\mathrm{tree}}
\le\frac3{32}
\le\frac18.
\]
Termwise,
\[
\frac{m^{m+1}}{m!}(4n\delta_{\mathrm L})^{m-2}
\le
\exp(1)^2\,m\,\zeta_{\mathrm{tree}}^{m-2}.
\]
The differentiated finite geometric sum yields, for \(n\ge1\),
\[
\begin{aligned}
\sum_{m=2}^n m\zeta_{\mathrm{tree}}^{m-2}
&=
\frac{
2-\zeta_{\mathrm{tree}}
-(n+1)\zeta_{\mathrm{tree}}^{n-1}
+n\zeta_{\mathrm{tree}}^n
}{(1-\zeta_{\mathrm{tree}})^2}\\
&\le
\frac{2-\zeta_{\mathrm{tree}}}
 {(1-\zeta_{\mathrm{tree}})^2}
\le3.
\end{aligned}
\]
For the first inequality, the removed tail equals
\(\zeta_{\mathrm{tree}}^{n-1}
 ((n+1)-n\zeta_{\mathrm{tree}})\ge0\).
For the second, \(1-\zeta_{\mathrm{tree}}\ge7/8\), and
\[
2-\zeta_{\mathrm{tree}}
\le2
\le3(7/8)^2
\le3(1-\zeta_{\mathrm{tree}})^2.
\]
Consequently,
\[
\sum_{m=2}^n
\frac{m^{m+1}}{m!}(4n\delta_{\mathrm L})^{m-2}
\le3\exp(1)^2
\le32.
\]
Substitution proves the required bound for the constructed coupling:
\[
\int d_{\mathrm H}(D,D')\,\gamma(dD,dD')
\le1024tn^2h_{\mathrm L}\delta_{\mathrm L}.
\]
% lean: CausalSmith.Stat.PrivateCateRoughdesign.tree_series_term_le
% lean: CausalSmith.Stat.PrivateCateRoughdesign.tree_series_sum_le_thirty_two
% lean: CausalSmith.Stat.PrivateCateRoughdesign.cubic_tree_envelope_sum_le
% lean: CausalSmith.Stat.PrivateCateRoughdesign.commonMassCoupling_hamming_cost_le_sparse

\item \textit{Transport and private contraction.}
For the remaining case \(n=0\), both datasets are empty, and
\[
d_{\mathrm H}(D,D')=0,
\qquad
\int d_{\mathrm H}(D,D')\,\gamma(dD,dD')
=0
=1024tn^2h_{\mathrm L}\delta_{\mathrm L}.
\]
Thus the coupling-cost bound holds over the full range of \(n\).

Since \(\gamma\) has the required marginals, the defining infimum of Hamming transport gives
\[
W_{\mathrm H}\!\left(Q_n,(P_0^{\mathrm{obs}})^{\otimes n}\right)
\le
\int d_{\mathrm H}(D,D')\,\gamma(dD,dD')
\le1024tn^2h_{\mathrm L}\delta_{\mathrm L}.
\]
Its marginals are probability laws, and the hypotheses on \(\epsilon\), the output space, and \(M\) permit application of
\cref{obj:lem:private-kernel-comparison}. That comparison, together with the defining infimum, supplies the coupling-cost inequality
\[
d_{\mathrm{TV}}\!\left(MQ_n,M(P_0^{\mathrm{obs}})^{\otimes n}\right)
\le
2\epsilon
\int d_{\mathrm H}(D,D')\,\gamma(dD,dD').
\]
Using the established cost bound and \(2\epsilon\ge0\), we conclude
\[
\begin{aligned}
d_{\mathrm{TV}}\!\left(MQ_n,M(P_0^{\mathrm{obs}})^{\otimes n}\right)
&\le
2\epsilon\bigl(1024tn^2h_{\mathrm L}\delta_{\mathrm L}\bigr)\\
&=
2048\epsilon tn^2h_{\mathrm L}\delta_{\mathrm L}.
\end{aligned}
\]
% lean: CausalSmith.Stat.PrivateCateRoughdesign.private_TV_le_coupling_cost
% lean: CausalSmith.Stat.PrivateCateRoughdesign.measurable_component_contraction
\end{enumerate}
\end{proof}
